\section{Results}
\label{sec:results}

\subsection{Validation against NEURON}
\label{sec:results-validation}

We validated jaxon against \texttt{pyfibers}-wrapped NEURON simulations
of all four axon models using the protocol of
Section~\ref{sec:methods-validation}.  Across 1\,152 distinct test
configurations spanning four models, 26 fiber diameters
($D \in [0.3,\,16.0]\,\mu\mathrm{m}$), eight stimulus waveforms
(monophasic cathodic and anodic, biphasic cathodic-first and
anodic-first, sinusoidal, sawtooth, exponentially-decaying and
Gaussian) and six pulse widths spanning $0.02$--$1\,\mathrm{ms}$,
\emph{100\,\% of configurations agree to within $1\%$ relative threshold
error vs.\ NEURON} (Table~\ref{tab:validation}).  The peak relative
error across all 432 MRG configurations is $0.39\%$, with median
$0.04\%$; the unmyelinated models perform similarly, with peak errors
of $0.46\%$ (Sundt, 240 configs), $0.17\%$ (Rattay, 240 configs) and
$0.02\%$ (Sweeney, 240 configs).  Conduction velocity agreement is
likewise sub-percent: median CV error rounds to zero for MRG, Sundt
and Sweeney over the manufacturer-supported diameter range, with
$0.24\%$ peak error in the smallest-diameter Sundt fiber.

For the MRG model, the agreement holds uniformly across the full
diameter range $D \in [5.7,\,16.0]\,\mu\mathrm{m}$ used in the
selectivity sweep of Section~\ref{sec:results-selectivity}: the
coupled $(V_i, V_{\mathrm{px}})$ backward-Euler solver introduces no
detectable drift from NEURON's analytic update at
$\Delta t = 5\,\mu\mathrm{s}$.  Figure~\ref{fig:validation-mrg} shows
the MRG strength--duration agreement at $D = 5.7\,\mu\mathrm{m}$ as an
exemplar; the analogous panels for the unmyelinated models are
included as Supplementary Figs.~S1--S3.

\begin{table}[h]
  \centering
  \caption{Strength--duration threshold agreement of jaxon vs.\
    \texttt{pyfibers}-wrapped NEURON, aggregated across all tested
    diameters, waveforms and pulse widths per model.  All four models
    agree to within $1\%$ on every single test configuration.}
  \label{tab:validation}
  \begin{tabular}{lrrrrr}
    \toprule
    Model    & $n_{\mathrm{cfg}}$ & median err.\ & 95\textsuperscript{th} \%ile & peak err.\ & \%\,$<1\%$ \\
    \midrule
    MRG      &  432 & 0.04\,\% & 0.20\,\% & 0.39\,\% & 100\,\% \\
    Sundt    &  240 & 0.34\,\% & 0.42\,\% & 0.46\,\% & 100\,\% \\
    Sweeney  &  240 & 0.03\,\% & 0.00\,\% & 0.02\,\% & 100\,\% \\
    Rattay   &  240 & 0.06\,\% & 0.14\,\% & 0.17\,\% & 100\,\% \\
    \bottomrule
  \end{tabular}
\end{table}

\begin{figure}[h]
  \centering
  \figmaybe[width=0.92\textwidth]{fig_mrg_analysis.png}
  \caption{Validation of jaxon against \texttt{pyfibers}-wrapped NEURON
    for the MRG myelinated fiber model.  Strength--duration thresholds
    are reproduced across the full clinically-relevant diameter range
    ($5.7$--$16\,\mu\mathrm{m}$) and across eight stimulus waveforms;
    every one of the 432 test configurations agrees to within $1\%$,
    with median error $0.04\%$.  Reproduction script:
    \texttt{experiments\_v2/validate\_mrg.py}.}
  \label{fig:validation-mrg}
\end{figure}

\subsection{Scaling on parallel fiber populations}
\label{sec:results-scaling}

The principal motivation for a JAX-based reimplementation is the
expectation that a GPU-vectorised forward pass will outperform the
CPU-parallel NEURON approach as the fiber count grows.
Figure~\ref{fig:scaling} reports the wall time per forward simulation
as a function of fiber count $N$ on a single NVIDIA A100 GPU, for
jaxon and \texttt{pyfibers} (the latter running 8-thread CPU parallel on a
matched Xeon CPU host).  We benchmarked all four primary models
(MRG, Sundt, Sweeney, Rattay) plus three additional models
(MRG with interpolated parameters, Schild94 and Schild97) for
comparison at $N \in \{1, 10, 100, 1\,000, 10\,000\}$.

At small $N$ the GPU launch overhead dominates and \texttt{pyfibers}
on CPU is faster: at $N=1$ jaxon takes $0.13$\,s vs.\
\texttt{pyfibers}'s $0.05$\,s for MRG (parity with pyfibers crosses at
$N\approx 10$).  Throughput then scales sub-linearly with $N$ on the
GPU because the vmap-batched forward pass amortises the per-step
launch cost across all fibers in the batch:
at $N=1\,000$ MRG, jaxon completes in $0.71\,\mathrm{s}$ vs.\
\texttt{pyfibers}'s $49.1\,\mathrm{s}$---a $69\times$ wall-time
reduction---and at $N=10\,000$ MRG, $8.3\,\mathrm{s}$ vs.\ $501\,\mathrm{s}$
($60\times$).  The unmyelinated Sweeney model, whose smaller per-fiber
state vector fits better in GPU registers, reaches a $122\times$
speedup at $N=10\,000$ ($5.0\,\mathrm{s}$ vs.\ $607\,\mathrm{s}$).
Table~\ref{tab:scaling} summarises the headline numbers.

The crossover point at which jaxon overtakes \texttt{pyfibers} is
$N\approx 10$ in every model we tested.  Standard fascicular
selectivity sweeps operate at $N\ge 200$ fibers per nerve realisation
\cite{hussain2024}, so jaxon is the right tool for every fascicular
problem larger than a single-fiber threshold lookup---and the
advantage grows monotonically with problem size.

\begin{table}[h]
  \centering
  \caption{Wall time per forward simulation (in seconds) and speedup
    of jaxon over \texttt{pyfibers} on representative fiber counts.
    jaxon runs on a single NVIDIA A100 GPU; pyfibers runs 8-thread
    CPU-parallel.  All timings exclude JAX compile time.}
  \label{tab:scaling}
  \begin{tabular}{lrrrrrr}
    \toprule
                            & \multicolumn{2}{c}{$N=100$}
                            & \multicolumn{2}{c}{$N=1\,000$}
                            & \multicolumn{2}{c}{$N=10\,000$} \\
    \cmidrule(lr){2-3}\cmidrule(lr){4-5}\cmidrule(lr){6-7}
    Model    & pyf.\ & speedup & pyf.\ & speedup & pyf.\ & speedup \\
    \midrule
    MRG      &   4.97 & $27\times$ & 49.1 & $69\times$ &  501 &  $60\times$ \\
    Sundt    &  17.55 & $30\times$ &  168 & $71\times$ & --- & --- \\
    Sweeney  &   6.49 & $37\times$ & 63.8 & $112\times$ & 607 & $122\times$ \\
    Rattay   &  16.04 & $24\times$ &  149 & $69\times$ & --- & --- \\
    \bottomrule
  \end{tabular}
\end{table}

\begin{figure}[h]
  \centering
  \figmaybe[width=0.85\textwidth]{fig_scaling.png}
  \caption{Wall time per forward simulation vs.\ fiber count $N$
    for jaxon on a single NVIDIA A100 GPU (lines) versus
    \texttt{pyfibers} (NEURON-wrapped, 8-thread CPU parallel) on a
    matched host (markers).  jaxon's curves are flat up to
    $N\approx 100$ as the GPU is under-utilised, then scale linearly.
    \texttt{pyfibers} scales linearly throughout.  The crossover is
    at $N\approx 10$ for every model; at $N=10\,000$ MRG, jaxon is
    $\sim 60\times$ faster than \texttt{pyfibers}.  Reproduction
    script: \texttt{experiments\_v2/scaling.py}.}
  \label{fig:scaling}
\end{figure}

\subsection{Selectivity optimisation on multi-fascicle nerves}
\label{sec:results-selectivity}

\subsubsection*{Single-fiber-diameter regime.}
We applied the rectangular-pulse and warm-started waveform optimisers
of Section~\ref{sec:methods-opt} to randomised Hussain-style
multi-fascicle nerve realisations
(Section~\ref{sec:methods-geom}) at a single fiber diameter of
$D = 5.7\,\mu\mathrm{m}$, a representative MRG fiber size used in the
prior literature~\cite{hussain2024}.  For each realisation we ran the
multi-start LBFGS rectangular optimiser to a budget of 30 inner
iterations per restart with 2 restarts, followed by 100 iterations of
the warm-started waveform Adam optimiser.

Across the seven seeds completed at the time of writing (of a planned
100-seed sweep currently running on the cluster), the rectangular
optimiser achieves a mean selectivity index of
$\mathrm{SI}=0.965$ (median $0.987$, range $0.900$--$1.000$);
$71\%$ of seeds reach $\mathrm{SI}\ge 0.95$ and $43\%$ reach the
theoretical maximum $\mathrm{SI}=1.000$ (all target fibers firing,
all off-target fibers silent).  No seed fell below $\mathrm{SI}=0.90$.
The warm-started waveform stage preserved the rectangular optimum
without regression on every seed where it was warm-started with the
LR\,=\,$5\times 10^{-4}$ schedule introduced in
Section~\ref{sec:methods-opt}.\needdata{Refresh from the 100-seed
sweep stats when the Phase 3 run on the a100 partition lands.}

Figure~\ref{fig:selectivity-example} shows a representative
cross-section for seed 0 ($\mathrm{SI}=0.987$), illustrating the
target-fascicle ``all fire'' / off-target-fascicle ``all silent''
activation pattern that LBFGS finds.  Figure~\ref{fig:selectivity-summary}
will summarise the distribution across all 100 seeds once the sweep
completes.

\begin{figure}[h]
  \centering
  \figmaybe[width=0.85\textwidth]{selectivity_violin_placeholder.pdf}
  \caption{Selectivity index distribution across $N=100$ Hussain-style
    multi-fascicle nerve realisations.  \textbf{a}, Violin plot of
    rectangular-pulse SI (left) and warm-started waveform SI (right);
    points show individual seeds, horizontal bars mark mean and IQR.
    \textbf{b}, Per-seed scatter of warm-started waveform SI vs
    rectangular SI showing that the waveform step preserves
    (refines but does not destabilise) the rectangular optimum.
    \needdata{Generate from analyze\_selectivity\_sweep.py after Phase 3
    completes.}}
  \label{fig:selectivity-summary}
\end{figure}

\begin{figure}[h]
  \centering
  \figmaybe[width=0.85\textwidth]{fig_seed_0000_rect_xsection.png}
  \caption{Example optimised activation pattern for seed 0.
    Nerve cross-section showing the eight fascicles (target = filled
    blue, off-target = hollow red), per-fiber activation
    (solid = fired, hollow = silent), and the six-contact cuff
    annotated by polarity (green = cathodic, orange = anodic) and
    amplitude.  All target fascicles fire; all off-target fascicles
    are silent, yielding $\mathrm{SI} = +1.000$ for this seed.}
  \label{fig:selectivity-example}
\end{figure}

\subsubsection*{Mixed-diameter regime.}
A second sweep targeted the harder problem of \emph{type
selectivity}---preferentially activating
small ($D = 5.7\,\mu\mathrm{m}$) target fibers without recruiting
co-located large ($D = 14.0\,\mu\mathrm{m}$) off-target
fibers---under the rationale that the differential chronaxie between
the two populations should be exploitable by waveform shaping.
Across all attempted regularisation strengths, soft-proxy
temperatures, and Adam learning rates, the optimiser invariably
converged to the all-fibers-fire equilibrium ($\mathrm{SI} = 0$
for both rectangular and waveform stages; see
Section~\ref{sec:discussion} for analysis of why the loss surface
favours this attractor under spatially homogeneous extracellular
forcing in our six-contact cuff geometry).  We report this negative
result rather than omit it: it delineates the regime in which the
present cuff geometry and loss formulation are sufficient
(single-diameter spatial selectivity) versus the regime in which an
additional design lever---most plausibly anatomically realistic
$V_e$ with target-fascicle-adjacent contacts---is required.

\subsection{Mixed-diameter limitation and outlook}
\label{sec:results-mixed-diam}

\needdata{Optional figure: the SI=0 plateau loss curves for the
mixed-diameter case, plus a chronaxie comparison panel between 5.7 and
14 \(\mu\)m MRG fibres showing how close the SD curves are at typical
PNS pulse widths.  This motivates the Discussion claim that waveform
shape alone is insufficient.}
